\documentclass[letterpaper,10pt]{article}

\usepackage[T1]{fontenc}
\usepackage{lmodern}
\usepackage[left=0.8in,right=0.8in,top=0.7in,bottom=0.7in]{geometry}
\usepackage{hyperref}
\usepackage{titlesec}
\usepackage{enumitem}
\usepackage{tabularx}
\usepackage{microtype}
\usepackage{needspace}

\hypersetup{colorlinks=true,urlcolor=blue,linkcolor=blue}
\urlstyle{same}
\pagestyle{empty}
\linespread{1.025}
\raggedbottom
\clubpenalty=10000
\widowpenalty=10000
\setlength{\parindent}{0pt}
\setlength{\parskip}{4pt}
\renewcommand{\arraystretch}{1.08}
\setlist[itemize]{leftmargin=1.4em,itemsep=4pt,topsep=5pt,parsep=0pt}
\titleformat{\section}{\large\bfseries}{}{0em}{}[\titlerule]
\titlespacing*{\section}{0pt}{12pt}{7pt}

\newcommand{\entry}[2]{%
  \noindent\begin{tabularx}{\textwidth}{@{}X r@{}}%
    \textbf{#1} & #2 \\
  \end{tabularx}\par\nopagebreak
}

\begin{document}

\begin{center}
    {\LARGE\textbf{Andrei Medvedev}}\\[5pt]
    Seattle, WA
    \quad $|$ \quad
    \href{mailto:am334@uw.edu}{am334@uw.edu}
    \quad $|$ \quad
    \href{https://github.com/mdvdvandrei}{github.com/mdvdvandrei}
\end{center}

\section{Summary}
University of Washington PhD student studying ocean circulation and climate variability and applying generative ML methods, mainly score-based diffusion models, to ocean and climate modeling, data assimilation, and probabilistic reconstruction from sparse observations.

\section{Education}

\entry{University of Washington, Seattle, WA}{2024 -- June 2029 (expected)}
\textit{PhD in Physical Oceanography}, GPA 3.95/4.00\\
Advisors: Georgy Manucharyan and Kyle Armour\\
\textit{Coursework:} Geophysical Fluid Dynamics I and II, Physics of Ocean Circulation, Advanced Physical Climatology

\smallskip
\entry{Lomonosov Moscow State University, Moscow, Russia}{2022 -- 2024}
\textit{MS in Physics, Atmospheric Physics}, GPA 3.92/4.00\\
\textit{Thesis:} Bayesian inference of 21st-century river runoff change from a multi-model climate ensemble

\smallskip
\entry{Lomonosov Moscow State University, Moscow, Russia}{2018 -- 2022}
\textit{BS in Physics, Applied Mathematics}\\
\textit{Thesis:} Deep reinforcement learning for an autonomous pursuit and interception task

\section{Research Experience}

\entry{Graduate Research Assistant, University of Washington \textnormal{(School of Oceanography)}}{2024 -- Present}

\textbf{Reconstructing the historical ocean with generative data assimilation} (first author; \href{https://doi.org/10.22541/essoar.15011822/v1}{preprint}, \href{https://github.com/mdvdvandrei/GenDA-historical-3d-ocean-reconstruction}{code})
\begin{itemize}
    \item Showed that historical observations determine the large-scale upper ocean back to 1950, including interior modes (Pacific thermocline tilt, Labrador Sea temperature at 1500~m), while the deep ocean, the Southern Ocean before the 1990s and Antarctic sea ice remain unconstrained.
    \item Produced a monthly 3D ocean reconstruction for 1950 to 2020 in which unassimilated sea surface height tracks satellite altimetry and basin-scale heat content agrees with observational products after 1990.
    \item Trained a video diffusion model (EDM, cBottle architecture, 73.5M parameters) from scratch on a 50-member climate-model ensemble and sampled posteriors via likelihood guidance (score-based data assimilation) from sparse observations, for all 852 months jointly.
\end{itemize}

\textbf{AMOC changes since 1900 from sea-surface observations} (first author, \textit{Nature Geoscience}, under review)
\begin{itemize}
    \item Found no statistically significant AMOC decline over the 20th century, in contrast to the widely used subpolar SST index, but a weakening over recent decades that is consistent with RAPID observations and smaller than in the climate model multi-model mean.
    \item Showed that the absence of a 20th-century decline is consistent with compensating greenhouse-gas and aerosol forcing, and that the recent weakening is statistically significant only in the non-Ekman component.
    \item Reconstructed AMOC at 26.5\textdegree N from basin-wide SST since the late 19th century and SSH since 1993 with deep learning trained on 18 CMIP6 models, giving yearly uncertainties and $\mathbf{r = 0.82}$ against RAPID.
\end{itemize}

\textbf{Additional projects}
\begin{itemize}
    \item Extended the NeurOST sea-level product (deep learning on satellite altimetry and SST) to a 2003 to 2025 record.
\end{itemize}

\smallskip
\needspace{6\baselineskip}
\entry{Research Assistant, A.M. Obukhov Institute of Atmospheric Physics}{2022 -- 2024}
\textit{Advisor: Alexey V. Eliseev}
\begin{itemize}
    \item Applied Bayesian inference to combine climate-model projections into probabilistic estimates of 21st-century river runoff change; published in \textit{Izvestiya, Atmospheric and Oceanic Physics} (2024).
\end{itemize}

\section{Professional Experience}

\entry{Machine Learning Intern (NLP), SkillCup \textnormal{(edtech startup)}}{2023}
\begin{itemize}
    \item Distilled a larger LLM into FRED-T5-large (820M parameters) for quiz generation from course texts and lecture transcripts, via synthetic pairs from the larger model plus a curated question bank, deployed with FastAPI.
\end{itemize}

\section{Technical Skills}

\textbf{Programming:} Python, C++, Bash, Linux, Git\\
\textbf{ML:} PyTorch (DDP), Hydra; diffusion and video diffusion (EDM), flow matching, diffusion bridges, guided posterior sampling, transformers (fine-tuning, distillation), reinforcement learning\\
\textbf{Statistics:} Bayesian inference, stochastic processes, time-series modeling, regression with correlated noise, uncertainty calibration, proper scoring rules (CRPS, energy score), ensemble methods\\
\textbf{Environmental data:} in-situ ocean profiles, satellite altimetry and SST, CMIP6 ensembles, HEALPix grids\\
\textbf{Data \& Systems:} NumPy, Pandas, xarray, Dask, Zarr, NetCDF; Docker, Podman, Apptainer; PBS/Slurm

\section{Publications \& Manuscripts}

\textbf{A. I. Medvedev}, H. Wei, K. C. Armour, K. Srinivasan, A. Solodoch, A. L. Stewart, and G. E. Manucharyan.
\textit{Basin-wide sea-surface observations indicate recent AMOC weakening without a 20th-century decline.}
Under review, \textit{Nature Geoscience}. Preprint on \href{https://doi.org/10.22541/essoar.15002298/v1}{ESS Open Archive}.

\smallskip
\textbf{A. I. Medvedev}, V. T. Cooper, and G. E. Manucharyan.
\textit{Reconstructing the global ocean from the historical observing record with generative data assimilation.}
Preprint on \href{https://doi.org/10.22541/essoar.15011822/v1}{ESS Open Archive}, to be submitted to \textit{JAMES}.

\smallskip
H. Wei, A. L. Stewart, \textbf{A. I. Medvedev}, K. Srinivasan, A. Solodoch, G. E. Manucharyan, and A. McC. Hogg.
\textit{Deep learning unlocks satellite-based monitoring of the global ocean overturning circulation.}
Preprint on \href{https://essopenarchive.org/doi/abs/10.22541/essoar.15004679}{ESS Open Archive}.

\smallskip
\textbf{A. I. Medvedev}, A. V. Eliseev, and I. I. Mokhov (2024).
\textit{Bayesian estimates of changes in Russian river runoff in the 21st century based on the CMIP6 ensemble model simulations.}
\textit{Izvestiya, Atmospheric and Oceanic Physics}.

\section{Presentations}

\entry{AGU Fall Meeting \textnormal{(two oral talks, generative data assimilation and AMOC reconstruction)}}{2026}
\entry{Graduate Climate Conference, Pack Forest, WA \textnormal{(poster, generative data assimilation)}}{2026}
\entry{AI4Ocean Workshop \textnormal{(lightning talk, score-based data assimilation)}}{2026}
\entry{AGU Fall Meeting \textnormal{(poster, machine learning reconstruction of historical AMOC variability)}}{2025}

\section{Awards \& Fellowships}

\entry{NASA FINESST Fellowship (\$150,000)}{2025 -- 2028}
Competitive federal graduate research fellowship (about 13\% selection rate)

\entry{Top-Off Award, University of Washington}{2024}
\entry{State Scholarship, Lomonosov MSU (top 10\% of students)}{2022 -- 2023}

\end{document}
